\documentclass[10pt,a4paper]{amsart}
\usepackage{amsmath,amssymb,amsthm,graphicx}
\usepackage{verbatim}
\newenvironment{reference}{\comment}{\endcomment}
%\newenvironment{reference}{}{}
\newenvironment{slogan}{\comment}{\endcomment}
\newenvironment{history}{\comment}{\endcomment}
% For commutative diagrams we use Xy-pic
\usepackage[all]{xy}

% We use 2cell for 2-commutative diagrams.
\xyoption{2cell}
% Theorem environments.
%
\theoremstyle{plain}
\newtheorem{theorem}[subsection]{Theorem}
\newtheorem{proposition}[subsection]{Proposition}
\newtheorem{lemma}[subsection]{Lemma}

\theoremstyle{definition}
\newtheorem{definition}[subsection]{Definition}
\newtheorem{example}[subsection]{Example}
\newtheorem{exercise}[subsection]{Exercise}
\newtheorem{situation}[subsection]{Situation}

\theoremstyle{remark}
\newtheorem{remark}[subsection]{Remark}
\newtheorem{remarks}[subsection]{Remarks}

\numberwithin{equation}{subsection}

% Macros
%
\def\lim{\mathop{\mathrm{lim}}\nolimits}
\def\colim{\mathop{\mathrm{colim}}\nolimits}
\def\Spec{\mathop{\mathrm{Spec}}}
\def\Hom{\mathop{\mathrm{Hom}}\nolimits}
\def\Ext{\mathop{\mathrm{Ext}}\nolimits}
\def\SheafHom{\mathop{\mathcal{H}\!\mathit{om}}\nolimits}
\def\SheafExt{\mathop{\mathcal{E}\!\mathit{xt}}\nolimits}
\def\Sch{\mathit{Sch}}
\def\Mor{\mathop{\mathrm{Mor}}\nolimits}
\def\Ob{\mathop{\mathrm{Ob}}\nolimits}
\def\Sh{\mathop{\mathit{Sh}}\nolimits}
\def\NL{\mathop{N\!L}\nolimits}
\def\CH{\mathop{\mathrm{CH}}\nolimits}
\def\proetale{{pro\text{-}\acute{e}tale}}
\def\etale{{\acute{e}tale}}
\def\QCoh{\mathit{QCoh}}
\def\Ker{\mathop{\mathrm{Ker}}}
\def\Im{\mathop{\mathrm{Im}}}
\def\Coker{\mathop{\mathrm{Coker}}}
\def\Coim{\mathop{\mathrm{Coim}}}

% Boxtimes
%
\DeclareMathSymbol{\boxtimes}{\mathbin}{AMSa}{"02}

%
% Macros for moduli stacks/spaces
%
\def\QCohstack{\mathcal{QC}\!\mathit{oh}}
\def\Cohstack{\mathcal{C}\!\mathit{oh}}
\def\Spacesstack{\mathcal{S}\!\mathit{paces}}
\def\Quotfunctor{\mathrm{Quot}}
\def\Hilbfunctor{\mathrm{Hilb}}
\def\Curvesstack{\mathcal{C}\!\mathit{urves}}
\def\Polarizedstack{\mathcal{P}\!\mathit{olarized}}
\def\Complexesstack{\mathcal{C}\!\mathit{omplexes}}
% \Pic is the operator that assigns to X its picard group, usage \Pic(X)
% \Picardstack_{X/B} denotes the Picard stack of X over B
% \Picardfunctor_{X/B} denotes the Picard functor of X over B
\def\Pic{\mathop{\mathrm{Pic}}\nolimits}
\def\Picardstack{\mathcal{P}\!\mathit{ic}}
\def\Picardfunctor{\mathrm{Pic}}
\def\Deformationcategory{\mathcal{D}\!\mathit{ef}}
\usepackage{cleveref}
\crefname{originalreference}{original reference}{original references}
\Crefname{originalreference}{Original reference}{Original references}
\makeatletter
\newlabel{more-algebra-lemma-base-change-RHom}{{more-algebra-lemma-base-change-RHom}{source}{}{sourcepointer-1}{}}
\newlabel{more-algebra-lemma-base-change-RHom@cref}{{[originalreference][0][]more-algebra-lemma-base-change-RHom}{[1][0][]source}{}{}{}}
\makeatother
\crefname{originalreference}{original reference}{original references}
\Crefname{originalreference}{Original reference}{Original references}
\makeatletter
\newlabel{more-algebra-lemma-finite-injective-dimension-Noetherian-local}{{more-algebra-lemma-finite-injective-dimension-Noetherian-local}{source}{}{sourcepointer-2}{}}
\newlabel{more-algebra-lemma-finite-injective-dimension-Noetherian-local@cref}{{[originalreference][0][]more-algebra-lemma-finite-injective-dimension-Noetherian-local}{[1][0][]source}{}{}{}}
\makeatother
\title{Base change of a normalized dualizing complex along a flat local map with the same maximal ideal}
\author{The Stacks Project Authors}
\date{stacks-project:dualizing.tex:local-snapshot-20261012; GNU FDL 1.2 or later}
\begin{document}
\maketitle
\noindent\textit{Editorial note.} All delivered mathematical text is the author's own, copied from the pinned Stacks Project chapter source sources/stacks-dualizing/pinned-original-source.tex; no statement or proof step was rewritten, reordered or removed.\par\medskip
\noindent\textit{Editorial note.} Preamble provenance: the Stacks Project keeps one single shared preamble for every chapter; this card takes it directly from the staged sources/stacks-dualizing/preamble.tex as three verbatim spans (lines 8-12, 14-18 and 154-217), so no derived file was written and nothing was added or redefined. The wrapper class and the base packages are supplied by the assembler.\par\medskip
\noindent\textit{Editorial note.} Bibliography provenance: neither the counted statement nor its proof contains a citation (this lane rescans the whole span before writing the card), so the card carries no bibliography and no entry was reconstructed.\par\medskip
\noindent\textit{Editorial note.} Selection provenance: this card was opened under the file-based control contract (CONTROL.md), with the mechanical pre-card gate run on its statement and proof interval before the card was written.\par\medskip
\noindent\textit{Allowed standard prerequisites (editorial).}\begin{itemize}
\item The base change statement for internal homs along flat maps. Invoked inside the counted proof (source lines 4266-4272 and 4277-4281) as a cross-chapter pointer. Transfers the internal hom computations along the base change.
\item The criterion for finite injective dimension of a complex over a Noetherian local ring with a single nonvanishing local cohomology module. Invoked inside the counted proof (source lines 4273-4276) as a cross-chapter pointer. Gives finite injective dimension of the base change.
\end{itemize}\medskip
\ifdefined\hypertarget\hypertarget{sourcepointer-1}{}\fi
\noindent\textit{Source reference.} Original reference more-algebra-lemma-base-change-RHom: sources/stacks-dualizing/pinned-original-source.tex, source lines 2719--2719. Navigation pointer only: the counted text refers to the Stacks item more-algebra-lemma-base-change-RHom, which belongs to another chapter of the project and is not part of this card. The pointer carries the target's own label as its number because this chapter's staged snapshot does not contain the other chapter's tag table; its content is not reproduced here.\par\medskip
\ifdefined\hypertarget\hypertarget{sourcepointer-2}{}\fi
\noindent\textit{Source reference.} Original reference more-algebra-lemma-finite-injective-dimension-Noetherian-local: sources/stacks-dualizing/pinned-original-source.tex, source lines 4074--4074. Navigation pointer only: the counted text refers to the Stacks item more-algebra-lemma-finite-injective-dimension-Noetherian-local, which belongs to another chapter of the project and is not part of this card. The pointer carries the target's own label as its number because this chapter's staged snapshot does not contain the other chapter's tag table; its content is not reproduced here.\par\medskip
\begin{lemma}
\label{lemma-flat-unramified}
Let $A \to B$ be a local homomorphism of Noetherian local rings.
Let $\omega_A^\bullet$ be a normalized dualizing complex.
If $A \to B$ is flat and $\mathfrak m_A B = \mathfrak m_B$,
then $\omega_A^\bullet \otimes_A B$ is a normalized dualizing
complex for $B$.
\end{lemma}

\begin{proof}
It is clear that $\omega_A^\bullet \otimes_A B$ is in $D^b_{\textit{Coh}}(B)$.
Let $\kappa_A$ and $\kappa_B$ be the residue fields of $A$ and $B$.
By More on Algebra, Lemma \ref{more-algebra-lemma-base-change-RHom}
we see that
$$
R\Hom_B(\kappa_B, \omega_A^\bullet \otimes_A B) =
R\Hom_A(\kappa_A, \omega_A^\bullet) \otimes_A B =
\kappa_A[0] \otimes_A B = \kappa_B[0]
$$
Thus $\omega_A^\bullet \otimes_A B$ has finite injective dimension by
More on Algebra, Lemma
\ref{more-algebra-lemma-finite-injective-dimension-Noetherian-local}.
Finally, we can use the same arguments to see that
$$
R\Hom_B(\omega_A^\bullet \otimes_A B, \omega_A^\bullet \otimes_A B) =
R\Hom_A(\omega_A^\bullet, \omega_A^\bullet) \otimes_A B = A \otimes_A B = B
$$
as desired.
\end{proof}

\end{document}
